\par\Needspace{6\baselineskip}\medskip
\begin{center}\hypertarget{native-196-article146}{}\pdfbookmark[2]{146.}{article-146}146.\end{center}
Zur Berechnung der Bahnelemente verwenden wir in der ersten Kombination $r',r'',2f$ und das korrigierte Zeitintervall zwischen der zweiten und dritten Beobachtung. Das Produkt dieses Intervalls mit der Größe $k$ (Art. 1) bezeichnen wir mit $\theta$. In der zweiten Kombination verwenden wir $r,r',2f''$ und das Zeitintervall zwischen der ersten und zweiten Beobachtung und setzen sein Produkt mit $k$ $=\theta''$. Die Rechnung nach dem in Art.$^{\prime}$ 88—105 dargelegten Verfahren ist jeweils nur bis zu der dort mit $y$ bezeichneten Größe fortzuführen. Ihren Wert bezeichnen wir in der ersten Kombination mit $\eta$, in der zweiten mit $\eta''$. Man setze dann
\[
\frac{\theta''\eta}{\theta\eta''}=P',\qquad
\frac{r'r'\theta\theta''}{rr''\eta\eta''\cos f\cos f'\cos f''}=Q',
\]
\SourceLeaf{197}Wenn die Werte der Größen $P,Q$, auf denen die ganze bisherige Rechnung beruht, selbst richtig waren, muss sich offenbar $P'=P,Q'=Q$ ergeben. Umgekehrt ist leicht einzusehen, dass für $P'=P,Q'=Q$ beide Rechnungen der Bahnelemente, bis zum Ende durchgeführt, genau dieselben Zahlen liefern. Diese stellen somit alle drei Beobachtungen exakt dar und lösen die Aufgabe vollständig. Ergibt sich hingegen nicht $P'=P,Q'=Q$, so nimmt man $P'-P,Q'-Q$ für $X$ und $Y$, sofern $P$ und $Q$ für $x$ und $y$ genommen wurden. Noch bequemer ist es, $\log P=x$, $\log Q=y$, $\log P'-\log P=X$, $\log Q'-\log Q=Y$ zu setzen.
Dann ist die Rechnung mit anderen Werten von $x,y$ zu wiederholen.
